% problem_id: S001-lemma-existence-plus-flat-base-change
% source_id: S001 (Stacks Project)
% source_file: stacks-morphisms.tex
% statement_locator: stacks-morphisms.tex lines 8671--8676
% proof_locator: stacks-morphisms.tex lines 8678--8771
% env: lemma
% pairing: tight (verified)
% status: complete (statement + author proof verbatim + reason + locators)
%
\begin{lemma}
\label{lemma-existence-plus-flat-base-change}
Let $f : \mathcal{X} \to \mathcal{Y}$ be a quasi-compact
morphism of algebraic stacks. Then formation of the scheme theoretic image
commutes with flat base change.
\end{lemma}

\begin{proof}
Let $\mathcal{Y}' \to \mathcal{Y}$ be a flat morphism of algebraic stacks.
Choose a scheme $V$ and a surjective smooth morphism $V \to \mathcal{Y}$.
Choose a scheme $V'$ and a
surjective smooth morphism $V' \to \mathcal{Y}' \times_\mathcal{Y} V$.
We may and do assume that $V = \coprod_{i \in I} V_i$ is a disjoint
union of affine schemes and that
$V' = \coprod_{i \in I} \coprod_{j \in J_i} V_{i, j}$
is a disjoint union of affine schemes with each $V_{i, j}$ mapping into $V_i$.
Let
\begin{enumerate}
\item $\mathcal{Z} \subset \mathcal{Y}$ be the scheme theoretic image of $f$,
\item $\mathcal{Z}' \subset \mathcal{Y}'$ be the scheme theoretic image
of the base change of $f$ by $\mathcal{Y}' \to \mathcal{Y}$,
\item $Z \subset V$ be the scheme theoretic image
of the base change of $f$ by $V \to \mathcal{Y}$,
\item $Z' \subset V'$ be the scheme theoretic image
of the base change of $f$ by $V' \to \mathcal{Y}$.
\end{enumerate}
If we can show that
(a) $Z = V \times_\mathcal{Y} \mathcal{Z}$,
(b) $Z' = V' \times_{\mathcal{Y}'} \mathcal{Z}'$, and
(c) $Z' = V' \times_V Z$
then the lemma follows: the inclusion
$\mathcal{Z}' \to \mathcal{Z} \times_\mathcal{Y} \mathcal{Y}'$
(Lemma \ref{lemma-factor-factor})
has to be an isomorphism because after base change by the surjective
smooth morphism $V' \to \mathcal{Y}'$ it is.

\medskip\noindent
Proof of (a). Set $R = V \times_\mathcal{Y} V$.
By Properties of Stacks, Lemma
\ref{stacks-properties-lemma-substacks-presentation}
the rule $\mathcal{Z} \mapsto \mathcal{Z} \times_\mathcal{Y} V$
defines a $1$-to-$1$ correspondence between closed substacks
of $\mathcal{Y}$ and $R$-invariant closed subspaces of $V$.
Moreover, $f : \mathcal{X} \to \mathcal{Y}$ factors through $\mathcal{Z}$
if and only if the base change
$g : \mathcal{X} \times_\mathcal{Y} V \to V$ factors through
$\mathcal{Z} \times_\mathcal{Y} V$.
We claim: the scheme theoretic image $Z \subset V$ of $g$
is $R$-invariant. The claim implies (a) by what we just said.

\medskip\noindent
For each $i$ the morphism $\mathcal{X} \times_\mathcal{Y} V_i \to V_i$
is quasi-compact and hence $\mathcal{X} \times_\mathcal{Y} V_i$ is
quasi-compact. Thus we can choose an affine scheme $W_i$ and a
surjective smooth morphism $W_i \to \mathcal{X} \times_\mathcal{Y} V_i$.
Observe that $W = \coprod W_i$ is a scheme endowed with
a smooth and surjective morphism $W \to \mathcal{X} \times_\mathcal{Y} V$
such that the composition $W \to V$ with $g$ is quasi-compact.
Let $Z \to V$ be the scheme theoretic image of $W \to V$, see
Morphisms, Section
\ref{morphisms-section-scheme-theoretic-image} and
Morphisms of Spaces, Section
\ref{spaces-morphisms-section-scheme-theoretic-image}.
It follows from Lemma \ref{lemma-cover-upstairs}
that $Z \subset V$ is the scheme theoretic image of $g$.
To show that $Z$ is $R$-invariant we claim that both
$$
\text{pr}_0^{-1}(Z), \text{pr}_1^{-1}(Z) \subset R = V \times_\mathcal{Y} V
$$
are the scheme theoretic image of $\mathcal{X} \times_\mathcal{Y} R \to R$.
Namely, we first use Morphisms of Spaces, Lemma
\ref{spaces-morphisms-lemma-flat-base-change-scheme-theoretic-image}
to see that $\text{pr}_0^{-1}(Z)$ is the scheme theoretic image
of the composition
$$
W \times_{V, \text{pr}_0} R = W \times_\mathcal{Y} V \to
\mathcal{X} \times_\mathcal{Y} R \to R
$$
Since the first arrow here is surjective and smooth we see that
$\text{pr}_0^{-1}(Z)$ is the scheme theoretic image of
$\mathcal{X} \times_\mathcal{Y} R \to R$. The same argument applies
that $\text{pr}_1^{-1}(Z)$. Hence $Z$ is $R$-invariant.

\medskip\noindent
Statement (b) is proved in exactly the same way as one proves (a).

\medskip\noindent
Proof of (c). Let $Z_i \subset V_i$ be the scheme theoretic image
of $\mathcal{X} \times_\mathcal{Y} V_i \to V_i$ and let
$Z_{i, j} \subset V_{i, j}$ be the scheme theoretic image of
$\mathcal{X} \times_\mathcal{Y} V_{i, j} \to V_{i, j}$.
Clearly it suffices to show that the inverse image of $Z_i$
in $V_{i, j}$ is $Z_{i, j}$. Above we've seen that
$Z_i$ is the scheme theoretic image of $W_i \to V_i$
and by the same token $Z_{i, j}$ is the scheme theoretic
image of $W_i \times_{V_i} V_{i, j} \to V_{i, j}$.
Hence the equality follows from the case of schemes
(Morphisms, Lemma
\ref{morphisms-lemma-flat-base-change-scheme-theoretic-image})
and the fact that $V_{i, j} \to V_i$ is flat.
\end{proof}
